% ============================================================
% Section 74: 一维双原子分子链的δ函数势与能隙
% ============================================================
\section{固体理论：一维双原子分子链的$\delta$函数势与能隙}
\label{sec:1d-diatomic-delta}

\begin{modelbox}[题目：一维双原子分子链的近自由电子能带]
考虑由 $N$ 个双原子分子组成的长度为 $L=Na$ 的一维固体。分子中的原子间距为 $b$（$b<a/2$），相邻分子的中心相距 $a$。势能为中心在每个原子上的 $\delta$函数之和：
\[
V(x) = -A\sum_{n=0}^{N-1}\Bigl[\delta\Bigl(x-na+\frac{b}{2}\Bigr)+\delta\Bigl(x-na-\frac{b}{2}\Bigr)\Bigr],\quad A>0.
\]
\begin{enumerate}
    \item 在自由电子近似下（$V=0$），推导电子波矢的允许值和归一化波函数；
    \item 将势能表示为傅里叶级数 $V(x)=\sum_q V_q e^{iqx}$，求 $q$ 的允许值和系数 $V_q$；
    \item 假定 $A$ 很小，用近自由电子近似证明对某些 $k$ 存在能隙，导出能隙的一般公式，并证明第一布里渊区顶部的能隙正比于 $\cos(\pi b/a)$；
    \item 导出第一布里渊区中的状态数。若每个原子有一个电子，该物质是导体还是绝缘体？
    \item 若 $b=a/2$，各结果出现什么新情况？给予简要解释。
\end{enumerate}
\end{modelbox}

\begin{tipbox}[考点快速分析]
\begin{itemize}
    \item \textbf{周期边界条件}：$k=2\pi n/L$，$n$ 为整数
    \item \textbf{傅里叶展开}：周期 $a$ 的函数，$q=m\cdot2\pi/a$
    \item \textbf{近自由电子能隙}：$E_g=2|V_G|$，$V_G$ 为对应倒格矢的傅里叶系数
    \item \textbf{导电性判据}：满带绝缘，部分填充为金属
    \item \textbf{$b=a/2$ 的特殊性}：双原子链退化为单原子链，布里渊区加倍
\end{itemize}
\end{tipbox}

\begin{derivationbox}[标准解题过程]
\textbf{Step 1: 自由电子波矢与归一化波函数}

$V(x)=0$ 时，定态薛定谔方程为
\begin{equation}
-\frac{\hbar^2}{2m}\frac{d^2\Psi}{dx^2}=E\Psi.
\end{equation}
解为平面波 $\Psi_k(x)=Ce^{ikx}$。采用周期性边界条件 $\Psi(0)=\Psi(L)$：
\begin{equation}
e^{ikL}=1\quad\Longrightarrow\quad k=\frac{2\pi}{L}n,\quad n=0,\pm1,\pm2,\ldots
\end{equation}

归一化条件 $\int_0^L|\Psi|^2dx=1$ 给出 $C=1/\sqrt{L}$。故
\begin{equation}
\boxed{\Psi_k(x)=\frac{1}{\sqrt{L}}e^{ikx}.}
\end{equation}

\textbf{Step 2: 势能的傅里叶级数展开}

$V(x)$ 以 $a$ 为周期，展开为
\begin{equation}
V(x)=\sum_{q}V_q e^{iqx},\qquad q=m\frac{2\pi}{a},\quad m=0,\pm1,\pm2,\ldots
\end{equation}
系数
\begin{equation}
V_q=\frac{1}{a}\int_{-a/2}^{a/2}V(x)e^{-iqx}\,dx.
\end{equation}

在一个周期内，$V(x)$ 包含两个 $\delta$ 函数，分别位于 $x=-b/2$ 和 $x=b/2$：
\begin{equation}
V(x)=-A\Bigl[\delta\Bigl(x+\frac{b}{2}\Bigr)+\delta\Bigl(x-\frac{b}{2}\Bigr)\Bigr].
\end{equation}
利用 $\int\delta(x-x_0)f(x)\,dx=f(x_0)$：
\begin{equation}
V_q=-\frac{A}{a}\bigl(e^{iqb/2}+e^{-iqb/2}\bigr)=-\frac{2A}{a}\cos\Bigl(\frac{qb}{2}\Bigr).
\end{equation}
代入 $q=m\cdot2\pi/a$：
\begin{equation}
\boxed{V_0=-\frac{2A}{a},\qquad V_m=-\frac{2A}{a}\cos\Bigl(\frac{m\pi b}{a}\Bigr)\;(m\neq0).}
\end{equation}

\textbf{Step 3: 近自由电子近似下的能隙}

将波函数按平面波展开 $\Psi(x)=\sum_k C_k e^{ikx}$，代入薛定谔方程得中心方程：
\begin{equation}
\bigl(E_k^{(0)}-E\bigr)C_k+\sum_q V_q C_{k-q}=0,\qquad E_k^{(0)}=\frac{\hbar^2k^2}{2m}.
\end{equation}

在一维情况下，布里渊区边界 $k=\pm m\pi/a$ 处发生简并：$E_k^{(0)}=E_{k-2m\pi/a}^{(0)}$。考虑 $k=m\pi/a$ 和 $k'=-m\pi/a$ 两个简并平面波，在近自由电子近似（简并微扰）下，久期方程为
\begin{equation}
\begin{vmatrix}
E_k^{(0)}+V_0-E & V_{2m\pi/a} \\[4pt]
V_{-2m\pi/a} & E_{k'}^{(0)}+V_0-E
\end{vmatrix}=0.
\end{equation}
解得
\begin{equation}
E=V_0+E_k^{(0)}\pm|V_{2m\pi/a}|.
\end{equation}

因此在 $k=m\pi/a$ 处打开能隙：
\begin{equation}
E_g(m)=2|V_{2m\pi/a}|=\frac{4A}{a}\Bigl|\cos\Bigl(\frac{m\pi b}{a}\Bigr)\Bigr|.
\end{equation}

第一布里渊区顶部对应 $m=1$（$k=\pi/a$）：
\begin{equation}
\boxed{E_g=\frac{4A}{a}\Bigl|\cos\frac{\pi b}{a}\Bigr|\propto\Bigl|\cos\frac{\pi b}{a}\Bigr|.}
\end{equation}

\textbf{Step 4: 第一布里渊区的状态数与导电性}

第一布里渊区范围：$-\pi/a<k\le\pi/a$。允许的波矢 $k=2\pi n/L=2\pi n/(Na)$，间隔 $\Delta k=2\pi/L=2\pi/(Na)$。

布里渊区内波矢数目
\begin{equation}
N_k=\frac{2\pi/a}{2\pi/(Na)}=N.
\end{equation}
计入自旋简并，总状态数为 $2N$。

晶体共有 $2N$ 个原子（$N$ 个双原子分子），每个原子贡献 1 个电子，总电子数为 $2N$。这些电子恰好填满第一布里渊区的所有状态。

\textbf{结论}：满带不导电，该材料为\textbf{绝缘体}（更准确地说是半导体，若带隙较小）。

\textbf{Step 5: $b=a/2$ 时的特殊情况}

当 $b=a/2$ 时，原子位置为 $na\pm a/4$，即等间距排列，相邻原子间距均为 $a/2$。晶格由双原子链退化为\textbf{单原子链}，新的晶格常数 $a'=a/2$。

\begin{itemize}
    \item \textbf{倒格矢}：周期变为 $a/2$，倒格矢基矢 $G'=2\pi/a'=4\pi/a$。
    \item \textbf{布里渊区}：扩大为 $-\pi/a'$ 到 $\pi/a'$，即 $-2\pi/a$ 到 $2\pi/a$，宽度加倍。
    \item \textbf{状态数}：新原胞数 $=L/a'=Na/(a/2)=2N$。布里渊区内波矢数 $=2N$（间隔仍为 $\Delta k=2\pi/L$）。计入自旋，总状态数 $=4N$。
    \item \textbf{电子填充}：总电子数仍为 $2N$。$2N$ 个电子只能填充 $4N$ 个状态的一半。
    \item \textbf{导电性}：部分填充的能带可导电，材料变为\textbf{导体}。
\end{itemize}

\textbf{物理本质}：$b=a/2$ 时晶格周期性增强（原胞缩小），布里渊区扩大，而电子总数不变，导致费米面落入布里渊区内部而非边界上，能带部分填充。
\end{derivationbox}

\begin{formulabox}[答案汇总]
\begin{center}
\begin{tabular}{ll}
\toprule
\textbf{问题} & \textbf{答案} \\
\midrule
(1) 波矢与波函数 & $k=2\pi n/L$；$\Psi_k(x)=L^{-1/2}e^{ikx}$ \\
(2) 傅里叶系数 & $V_0=-2A/a$；$V_m=-\frac{2A}{a}\cos\bigl(\frac{m\pi b}{a}\bigr)$ $(m\neq0)$ \\
(3) 能隙 & $E_g(m)=\frac{4A}{a}|\cos(m\pi b/a)|$；$m=1$ 时 $\propto|\cos(\pi b/a)|$ \\
(4) 状态数与导电性 & $2N$ 个状态；$2N$ 个电子填满；\textbf{绝缘体} \\
(5) $b=a/2$ & 单原子链，布里渊区加倍，电子部分填充；\textbf{导体} \\
\bottomrule
\end{tabular}
\end{center}
\end{formulabox}

\begin{insightbox}[物理图像]
\begin{itemize}
    \item 能隙大小直接由势场的傅里叶分量决定。$b\to0$ 时两原子靠近，$\cos(\pi b/a)\to1$，能隙最大；$b\to a/2$ 时 $\cos(\pi b/a)\to0$，能隙消失——这对应于晶格对称性提高、两种原子等价的情形。
    \item 本题是\textbf{佩尔斯相变}（Peierls transition）的一维原型：一维金属中的晶格畸变可以打开能隙，使系统从金属变为绝缘体。
    \item $b=a/2$ 时的金属--绝缘体转变说明：晶格结构的微小变化（原子相对位置）可以根本改变电子结构的拓扑和材料的导电性质。
\end{itemize}
\end{insightbox}
